\section{模型评价与推广}

本文从统一方程、守恒离散、连续事件和数值检验四个层面评价模型。四个子问题沿同一径向热—水分传递框架递进：问题一建立固定尺寸的预热基线，问题二引入状态相关物性，问题三将全域含水率阈值转化为连续事件，问题四在材料坐标中处理给定收缩轨迹并分离几何与物性作用。该结构保证了不同问题使用一致的物理量和数值口径。

模型的优点在于保留圆柱几何权重、界面通量连续性和全域阈值判据。环形有限体积法使相邻控制体共享同一界面通量；水分通量在 $K(C)$ 空间积分以处理扩散率非线性；表面采用半控制体 Robin 条件并求解真实表面状态。问题三和问题四在连续重构场上定位最大含水率事件，避免以展示位置或单元中心代替全域判据。问题四的四组合对照进一步区分收缩几何、整组物性及交互作用。

数值检验覆盖空间网格、时间积分、守恒关系、Robin 边界、连续场重构和独立基准。问题一在 $2048\to4096$ 网格加密后最大温度变化为 $1.504\times10^{-7}\,\mathrm K$，独立圆柱 Robin 热基准的最大偏差为 $1.804\times10^{-7}\,\mathrm K$；问题二同级加密后的最大温度和含水率差分别为 $1.163\times10^{-6}\,\mathrm K$ 和 $4.243\times10^{-7}\,\mathrm{kg/kg}$。问题三、四的 $N=512\to1024$ 连续达标时刻变化分别为 $0.132\,\mathrm s$ 和 $0.136\,\mathrm s$，干基水量收支残差均为 $10^{-14}$ 量级。上述结果支持当前离散与结果提取的一致性，不构成真实药材的实验误差上界。

敏感性对照表明，长期边界输入比网格误差更值得关注。将 $4$ h 后环境由末小时均值改为最后观测值，问题三和问题四的达标时刻分别提前约 $18.28$ min 和 $16.07$ min；问题四采用 PCHIP 半径插值时事件提前约 $12.75$ s。因此，长期预测必须注明环境延拓和半径插值口径。四组合结果显示，相对固定半径—附录 3 基准，几何效应为 $-32.230058$ h，物性效应为 $+72.374341$ h，交互效应为 $-46.526182$ h；这些数值是给定半径轨迹下的反事实机制对照，不等同于因果识别。

模型仍有明确边界：有效表面平衡浓度和附录经验物性尚未由本题实验反演；热方程未显式计入蒸发潜热；一维径向模型忽略端部传递；问题四采用均匀收缩及附件给定半径轨迹，未建立含水率到收缩率的反馈；$4$ h 后环境采用末小时均值延拓。因而本文结论适用于题面给定尺寸、边界轨迹、物性和收缩情景下的比较与预测，不应直接解释为所有药材和工艺条件下的普适规律。

若补充水活度、空气湿度、潜热、端部边界、收缩实验和独立剖面数据，可将模型扩展为空气—固体气固热质耦合的二维轴对称或三维模型，并对物性、边界和收缩参数开展不确定性传播。任何阶段降阶或参数敏感度分析都应以当前完整耦合解为参考，同时报告场误差、达标时间误差和计算成本。
